\chapter{Вещества, которые не задают тип нейрона}
\label{sec:othermediators}

В главе~\ref{sec:neurontypes} мы рассортировали нейроны по главному медиатору и получили десять типов (таблица~\ref{tab:types}). Теперь пора добавить осложнения. Нейрон выбрасывает не только главный медиатор, а в мозге кроме главных медиаторов работают и другие вещества. Они меняют то, как сеть передаёт и хранит сигналы.

Кроме адресной и широковещательной шин (раздел~\ref{sec:buses}) есть и третья: \emph{обратный канал}. Несколько веществ выбрасывает не отправитель, а \emph{получатель}, и они идут через щель назад, к выходу отправителя, и меняют, сколько медиатора тот выбросит в следующий раз. В сетях связи это похоже на подтверждение приёма, которое отправитель использует, чтобы настроить свою передачу. Именно этим каналом нейроны пользуются, когда меняют веса своих связей: по нему выход отправителя узнаёт об активности получателя, то есть о втором множителе правил обучения главы~\ref{sec:dofa}.

Полный список известных медиаторов длинен: больше сотни веществ, большинство из них --- короткие белковые цепочки, нейропептиды. Но все они укладываются в несколько классов, и внутри каждого класса вещества ведут себя сходно. Вещества, которые не бывают главным медиатором, собраны в таблице~\ref{tab:othermediators}, с теми же тремя вопросами, которые задал бы инженер: по какой шине идёт сигнал, быстро ли он действует и какой у него обычно знак. Тип нейрона они не задают: они выбрасываются вместе с главным медиатором, или выбрасываются не нейронами, или идут в обратную сторону.

\begin{table}[t]
\centering
\footnotesize
\setlength{\tabcolsep}{4pt}
\renewcommand{\arraystretch}{1.12}
\begin{tabular}{>{\raggedright}p{2.25cm}>{\raggedright}p{2.8cm}>{\raggedright}p{1.8cm}>{\raggedright}p{1.5cm}>{\centering}p{0.8cm}>{\raggedright\arraybackslash}p{4.3cm}}
\toprule
Класс & Вещество & Шина & Скорость & Знак & Роль в вычислении \\
\midrule
Аминокислоты & D-серин & локальная & медленно & И & второй ключ для рецептора глутамата: условие «И» \\
\midrule
Моноамины & следовые амины & широковещ. & медленно & $\pm$ & тонкая настройка моноаминовых каналов \\
\midrule
\multirow{2}{2.25cm}{Пурины}
 & АТФ & адресная & быстро и медленно & $+$ & сопутствующий медиатор; сигнал между нейронами и глией \\
 & аденозин & объёмная & медленно & $-$ & накапливается при активности: счётчик нагрузки~\cite{Dunwiddie2001} \\
\midrule
\multirow{8}{2.25cm}{Нейропептиды (больше сотни)}
 & энкефалины, эндорфины, динорфин & объёмная & медленно & $-$ & подавление передачи \\
 & вещество P & объёмная & медленно & $+$ & медленное усиление \\
 & нейропептид Y & объёмная & медленно & $-$ & медленное торможение \\
 & соматостатин & объёмная & медленно & $-$ & медленное торможение, сопутствует ГАМК \\
 & вазоактивный интестинальный пептид (ВИП) & объёмная & медленно & $+$ & растормаживание: торможение тормозящих нейронов \\
 & холецистокинин & объёмная & медленно & $\pm$ & сопутствует ГАМК в части тормозящих нейронов \\
 & орексин & широковещ. & медленно & $+$ & стабилизатор режима бодрствования \\
 & окситоцин, вазопрессин & широковещ. & медленно & $\pm$ & глобальные настройки поведения \\
\midrule
\multirow{2}{2.25cm}{Газы}
 & оксид азота NO & обратная, объёмная & медленно & $\pm$ & проходит сквозь мембраны; обратный сигнал при изменении весов~\cite{Garthwaite2008} \\
 & CO, H$_2$S & объёмная & медленно & $\pm$ & роль изучена мало \\
\midrule
Липиды & эндоканнабиноиды (анандамид, 2-АГ) & обратная & медленно & $-$ & получатель уменьшает выброс у отправителя: обратная связь по весу~\cite{WilsonNicoll2001} \\
\bottomrule
\end{tabular}
\caption{Вещества, которые не задают тип нейрона. \emph{Шина}, \emph{скорость} и \emph{знак} --- как в таблице~\ref{tab:types}; кроме того, объёмная шина --- медиатор расплывается в межклеточном пространстве вокруг места выброса, обратная --- идёт от получателя назад к отправителю, локальная --- действует рядом с местом выброса. «И» --- разрешающий сигнал: сам по себе ток не вызывает, но без него не срабатывает другой вход, как второй вход логического элемента «И». Нейропептиды почти всегда выбрасываются вместе с главным медиатором и в основном при высокой частоте импульсов~\cite{vandenPol2012}.}
\label{tab:othermediators}
\end{table}
